\section{Conclusion}
This paper proposed Q-band ART, an adaptive frequency-band suppression method for anti-rollover tracking of semi-trailer tank trucks. Instead of using liquid sloshing states or wheel-load-based LTR as online controller states, it suppresses predicted lateral-acceleration energy in critical vehicle-liquid bands. A reduced roll-sloshing model provides the frequency-domain safety interpretation, and the Q-band term is integrated with acceleration smoothing, soft constraints, preview-adaptive scheduling, ESO compensation, and hitch-angle-based parameter scheduling.

Continuous five-pins co-simulation and independent-DLC ablation show that Q-band ART avoids rollover where pure preview tracking and acceleration-amplitude limiting fail. Compared with full-state active steering, it provides comparable anti-rollover tracking with lower dependence on liquid-state observation, trailer-side feedback, and high-dimensional online models. The results indicate that input-side critical-band suppression is a practical complement to explicit rollover-state constraints. Future work will refine critical-band calibration across filling rates and integrate Q-band scheduling with combined steering and braking control.
